\section{Same-residual divisor compression}
\label{sec:anchor-compression}

The equal-residual edge is much more rigid than a generic pair of
cofactors.

\subsection{Primitive row-divisor coordinates}

Write
\begin{equation}
 g:=(d_m,d_n),\qquad
 d_m=ga,\qquad d_n=gb.
 \label{eq:row-divisor-coordinates}
\end{equation}
Then \((a,b)=1\).  Since \(d_m,d_n\) are squarefree, the three integers
\(g,a,b\) are squarefree and pairwise coprime.

\begin{theorem}[Exact anchor compression]
\label{thm:anchor-compression}
Under \eqref{eq:same-residual-anchor} and the squarefree support
\eqref{eq:squarefree-residuals}, there is a unique positive integer \(t\)
such that
\begin{equation}
 \boxed{R_m=bt,\qquad R_n=at.}
 \label{eq:anchor-cofactor-compression}
\end{equation}
Moreover \(t\) is squarefree, is coprime to \(gab\), and
\begin{equation}
 \boxed{t\mid |H|.}
 \label{eq:t-divides-H}
\end{equation}
The two largest-prime labels obey the exact line
\begin{equation}
 \boxed{nbP_m-maP_n=\frac{H}{t}.}
 \label{eq:anchor-prime-line}
\end{equation}
\end{theorem}

\begin{proof}
The residual equality and \eqref{eq:row-divisor-coordinates} give
\[
 aR_m=bR_n.
\]
Because \((a,b)=1\), one has \(b\mid R_m\) and \(a\mid R_n\).  Write
\[
 R_m=bt_m,\qquad R_n=at_n.
\]
Substitution gives \(t_m=t_n\), proving
\eqref{eq:anchor-cofactor-compression} and uniqueness.

The common residual is
\[
 d_mR_m=d_nR_n=gabt.
\]
It is squarefree by \eqref{eq:squarefree-residuals}.  Therefore \(t\) is
squarefree and coprime to \(gab\).  Substituting
\eqref{eq:anchor-cofactor-compression} into
\eqref{eq:anchor-prime-determinant} gives
\[
 t(nbP_m-maP_n)=H.
\]
This proves \eqref{eq:t-divides-H} and
\eqref{eq:anchor-prime-line}.
\end{proof}

\begin{corollary}[Divisor-many common residual cores]
\label{cor:divisor-many-anchor-cores}
For fixed \(m,n,h_0,d_m,d_n\), the scale \(t\) of an actual
same-residual anchor belongs to the finite set
\[
 \{t\in\N:t\mid |h_0(n-m)|\}.
\]
Hence the number of possible \(t\)-values is at most
\(\tau(|h_0(n-m)|)\).  If a physical cofactor window forces
\(t>|H|\), that row pair has no same-residual anchor in the window.
\end{corollary}

\begin{remark}[What is and is not compressed]
\label{rem:anchor-compression-scope}
The parameter \(t\) records the entire common residual after the fixed
row factors \(g,a,b\) are removed.  The theorem does not remove the
largest-prime constraints, the native-key labels, or the literal
coefficient weights.  In particular, \(\tau(|H|)\) is a count of
arithmetic core types, not a bound for the number or total mass of
physical rectangles.
\end{remark}

\subsection{A first cross-determinant consequence}

The same-residual substitution also rewrites the cross determinant as
\begin{equation}
 tq\bigl(bvP_mQ_n-auQ_mP_n\bigr)=H\Delta
 \label{eq:primitive-cross-determinant-preview}
\end{equation}
once the erasure coordinates \(q,u,v\) of
\cref{sec:erasure-coordinates} are introduced.  Thus every physical
rectangle later satisfies \(tq\mid H\Delta\).  This compatibility is
automatic for actual targets, but it does not imply \(tq\mid H\) and it
does not determine a M\"obius sign.
